\documentclass[12 pt, a4paper]{article}% тип документа, размер шрифта
\usepackage{cmap}	
\usepackage{hyperref}
\hypersetup{
    colorlinks=true,
    linkcolor=blue,
    urlcolor=blue,
}
\usepackage{mathtext}
\usepackage[T2A]{fontenc}%поддержка кириллицы в ЛаТеХ
\usepackage[utf8]{inputenc}%кодировка
\usepackage[english,russian]{babel}
\usepackage{indentfirst}

\usepackage{amsmath,amsfonts,amssymb,amsthm,mathtools} % AMS
\usepackage{amsmath}%удобная вёрстка многострочных формул, масштабирующийся текст в формулах, формулы в рамках и др.
\usepackage{amsfonts}%поддержка ажурного и готического шрифтов — например, для записи символа {\displaystyle \mathbb {R} } \mathbb {R} 
\usepackage{amssymb}%amsfonts + несколько сотен дополнительных математических символов
\frenchspacing%запрет длинного пробела после точки
\usepackage{setspace}%возможность установки межстрочного интервала
\usepackage{indentfirst}%пакет позволяет делать в первом абзаце после заголовка абзацный отступ
\onehalfspacing%установка полуторного интервала по умолчанию
\usepackage{graphicx}%подключение рисунков
\graphicspath{{images/}}%путь ко всем рисункам
\usepackage{caption}
\usepackage{float}%плавающие картинки
\usepackage{tikz} % это для чудо-миллиметровки
\usepackage{pgfplots}%для построения графиков
\pgfplotsset{compat=newest, y label style={rotate=-90},  width=10 cm}%версия пакета построения графиков, ширина графиков
\usepackage{pgfplotstable}%простое рисование табличек
\usepackage{lastpage}%пакет нумерации страниц
\usepackage{comment}%возможность вставлять большие комменты
\usepackage{python}%питончик
\usepackage{float}
%%%%% ПОЛЯ
\setlength\parindent{0pt} 
\usepackage[top = 2 cm, bottom = 2 cm, left = 1.5 cm, right = 1.5 cm]{geometry}
\setlength\parindent{0pt}
%%%%% КОЛОНТИТУЛЫ
\usepackage[shortlabels]{enumitem}

\usepackage{fancybox, fancyhdr}
\pagestyle{fancy} 
\fancyhead[L]{Легких Никита Б05-301}
\fancyhead[C]{2.5.1}
% ЛЮБАЯ ДОПОЛНИТЕЛЬНАЯ ИНФОРМАЦИЯ
\fancyfoot[R]{ФПМИ МФТИ}
\renewcommand{\footrulewidth}{0.3 mm}
 

\usepackage{graphicx}
\graphicspath{{pictures/}}
\DeclareGraphicsExtensions{.pdf,.png,.jpg}
\usepackage{wrapfig}
%%% Заголовок
%%% Новые команды
\newcommand{\z}[2]{{{\vspace{0.6cm} \textbf{{Задача {#1}}} \textit{{#2}\\ }}} }
\newcommand{\otv}{{\vspace{0.3cm} \textbf{Ответ: }}}
\newcommand{\uk}{\underline{\textit{Указание.}} }
\newcommand{\prim}{\underline{\textit{Примечание.}} }
\newcommand{\sol}[1]{{{\vspace{0.3cm} \textbf{{Задача {#1}} }\\ }}}



\begin{document} % конец преамбулы, начало документа

\begin{center}
\section*{Работа 2.5.1}
\section*{Измерение коэффициента поверхностного натяжения жидкости}
\end{center}
\textbf{Цель работы:} 1) измерение коэффициента поверхностного натя- жения исследуемой жидкости при разной температуре с использованием известного коэффициента поверхностного натяжения другой жидкости; 2) определение полной поверхностной энергии и теплоты, необходимой для изотермического образования единицы поверхности жидкости \\ \\
\textbf{В работе используются:} прибор Ребиндера с термостатом; исследуемые жидкости; стаканы.

\section*{Теоретическая справка}
Из-за наличия поверхностного слоя жидкости и ее искревления появляется избыточное давление которое можно найти по формула Лапласа: $$\Delta P = P_{inside} - P_{outside} = \sigma(\frac{1}{R_1} + \frac{1}{R_2}) = \frac{4\sigma}{D}$$ 
Данная формула будет лежаьт в основе нашего метода.

\section*{Установка}
Исследуемая жидкость (вода) наливается в сосуд $B$. Известная жидкость (спирт) наливается в сосуд $E$. Сосуды закрыты пробками. Через пробку сосуда, в котором проводятся измерения, проходит полая металлическая игла $С$, нижний конец которой погружен в жидкость, а верхний открыт в атмосферу. Если другой сосуд герметично закрыт, то в сосуде с иглой создается разрежение, и пузырьки воздуха начинают пробулькивать через жидкость. Поверхностное натяжение можно найти по величине разрежения, необходимого для прохождения пузырьков. При приоткрытом кране $К_1$ из аспиратора A по каплям вытекает вода, создавая разрежение, которое измеряется наклонным спиртовым манометром $М$. Показания манометра, умноженные на зависящий от наклона коэффициент (в нашем случае 0,2), дают давление в кгс/м$^2$ (1 кгс/м$^2$ = 9,8 Па). Чтобы пополнить запас воды, достаточно при помощи крана $К_2$ соединить нижнюю часть аспиратора с атмосферой и предварительно заполненной водой верхней частью.
Для стабилизации температуры исследуемой жидкости через рубашку $D$ непрерывно прогоняется вода из термостата.
Обычно кончик иглы лишь касается поверхности жидкости, чтобы исключить влияние гидростатического давления столба жидкости. Однако при измерении температурной зависимости коэффициента поверхностного натяжения возникает ряд сложностей. Во-первых, большая теплопроводность металлической трубки приводит к
тому, что температура на конце трубки заметно ниже, чем в глубине
жидкости. Во-вторых, тепловое расширение поднимает уровень жидкости при увеличении температуры. Это гидростатическое давление вычитается из падения лапласова давления вследствие уменьшения $\sigma$, и в опыте с анилином, например, наблюдаемый эффект меняет знак при высоте столба жидкости порядка пяти сантиметров.

Обе погрешности устраняем путем погружения иглы наполовину возможной глубины, при этом измеряемое давление $P = \Delta P +\rho g h$, но теперь нам надо определить величину $\rho g h$, для этого будем измерять $P_1 = \Delta P'$- давление около поверхности воды и $P_2 = \Delta P'' + \rho g h$. Из-за несжимаемости жидкости можно положить $\Delta P' = \Delta P''$, тогда отсюда $\rho g h = P_2 - P_1$
\begin{figure}[h]
\centering
\includegraphics[width=0.8\linewidth]{2.5.1/рис1.png}
\label{fig:mpr}
\end{figure}
\section*{Задание}
\begin{enumerate}
    \item В исправности установки убежден, пробулькивание происходит и в дальнейщем для его обеспечения всегда будет выставляться такая скорость образования капли в $K_1$, что будет образовываться 1 пузырь в 5 секунд
    \item Для выполнения условия первого пункта во всех случаях скорость падения капли- 1 в 5 секунд.
    \item Для нахождения максимального $\Delta P_{спирт}$ измерим несоклько различных значений $h$ снимаемых с микроманометра и воспользуемся формулой $P = 0.2 \cdot 9.81 \cdot h$
    \begin{center}
    \begin{tabular}{|l|l|l|l|l|l|}
    \hline
    $N$                    & 1    & 2    & 3    & 4    & 5    \\ \hline
    $h$, мм                & 40   & 41   & 40   & 39   & 40   \\ \hline
    $\Delta P_{спирт}$, Па & 78,5 & 80,4 & 78,5 & 76,5 & 78,5 \\ \hline
    \end{tabular}
    \end{center}
    Из таблицы $\rightarrow \Delta P_{спирт}= (78.5 \pm 1.9) \;Па$ \newpage
    Табличное значения коефициента поверхностного натяжения для спирта $\sigma = 22.75\cdot10^{-3}\; Н/м$, подставим все в формулу $R = \frac{2\sigma}{\Delta P_{спирт}} = 0.58\; мм$. Погрешность: $\delta R = R \cdot \frac{\delta(\Delta P_{спирт})}{\Delta P_{спирт}} = 0.01\; мм$\\
    При прямом измерении получаем что $D_{прямое} = (1.1 \pm 0.05)\; мм$, как видно с помощью $\Delta P_{спирт}$ получается сильно точнее, будем использовать значение $R$.
    \item Опускаем иглу в воду на такую глубину, чтобы выходящий пузырек не касался дна, измерим эту глубину двумя способами: \\
    \textbf{С помощью давления}, сначала игла только касается поверхности воды:
    \begin{center}
    \begin{tabular}{|l|r|r|r|r|r|}
    \hline
    $N$       & \multicolumn{1}{l|}{1} & \multicolumn{1}{l|}{2} & \multicolumn{1}{l|}{3} & \multicolumn{1}{l|}{4} & \multicolumn{1}{l|}{5} \\ \hline
    $h$, мм   & 114                    & 112                    & 111                    & 113                    & 111                    \\ \hline
    $P$, Па & 223.7                  & 219.7                  & 217.8                  & 221.7                  & 217.8                  \\ \hline
    \end{tabular}
    \end{center}
    Отсюда: $P_1 = (220.1 \pm 6.5) \;Па$, теперь же погружаем иглу: 
    \begin{center}
    \begin{tabular}{|c|c|c|c|c|c|}
    \hline
    $N$       & 1     & 2     & 3     & 4     & 5     \\ \hline
    $h$, мм   & 184   & 183   & 182   & 182   & 185   \\ \hline
    $P_2$, Па & 361.0 & 359.0 & 357.1 & 357.1 & 363.0 \\ \hline
    \end{tabular}
    \end{center}
    Воспользуемся очевидной формулой: $\Delta h = \Delta P'/ \rho g$, где $\rho = 1\; г/см^3,\;\; \Delta P' = P_2 - P_1$, отсюда $\Delta h = (14.2 \pm 0.2)\; мм$ \\
    \item \textbf{Прямое измерение}: $h_1 = (19.5\pm0.5)\; мм$, $h_2 = (6.5\pm0.5)\; мм$, тогда $h_2 - h_1 = (13\pm1)\; мм$ видно что значение $\Delta h$ попадает в область возможных значений $h_2 - h_1$, поэтому имеет право на жизнь и впоследствии будем пользоваться именно $\Delta h$ из-за большей точности $\Rightarrow \Delta P = \rho g \Delta h = (139.3\pm2.0)\; Па$.
    \item Снимаем требуемую зависимость $\sigma (T)$, идем с шагом $5^oC$:
    \begin{center}
    \begin{tabular}{|c|c|c|c|c|c|c|c|c|}
    \hline
    $T$, $^oC$ & $T$, K & $h_1$, мм & $h_2$, мм & $h_3$, мм & $<P>$, Па & $\delta <P>$, Па & $\sigma$, мН/м & $\delta \sigma$, мН/м \\ \hline
    21         & 294    & 184       & 183       & 182       & 219.74    & 3.96             & 63.7           & 2.2                   \\ \hline
    25.1       & 298.1  & 181       & 181       & 181       & 215.82    & 3.96             & 62.6           & 2.2                   \\ \hline
    30.1       & 303.1  & 178       & 180       & 179       & 211.90    & 3.96             & 61.4           & 2.2                   \\ \hline
    35.1       & 308.1  & 175       & 176       & 178       & 206.66    & 3.96             & 59.9           & 2.2                   \\ \hline
    41         & 314    & 174       & 173       & 174       & 201.43    & 3.96             & 58.4           & 2.2                   \\ \hline
    45.1       & 318.1  & 172       & 172       & 172       & 198.16    & 3.96             & 57.5           & 2.1                   \\ \hline
    50.1       & 323.1  & 170       & 169       & 170       & 193.58    & 3.96             & 56.1           & 2.1                   \\ \hline
    55         & 328    & 169       & 167       & 167       & 189.66    & 3.96             & 55.0           & 2.1                   \\ \hline
    58.2       & 331.2  & 167       & 167       & 166       & 187.70    & 3.96             & 54.4           & 2.1                   \\ \hline
    \end{tabular}
    \end{center}
    пересчет столбцов происходит с помощью следующих формул: \\
    $$<P> = \frac{h_1 + h_2 + h_3}{3}\cdot 0.2 \cdot 9.81 - \Delta P$$ 
    $$\sigma = \frac{<P>\cdot R}{2}$$ $$\delta(<P>) = \frac{\delta h}{3}\cdot 0.2 \cdot 9.81 + \delta(\Delta P)$$
    $$\delta \sigma = \sigma(\epsilon_{<P>} + \epsilon_R) = \sigma(\frac{\delta(<P>)}{<P>} + \frac{\delta R}{R})$$
    \item см прошлый пункт
    \item Строим график $\sigma(T)$ \begin{figure}[h]
    \centering
    \includegraphics[width=0.7\linewidth]{2.5.1/график сигма от T.png}
    \label{fig:mpr} 
    \end{figure} \\
    Из данного графика: $$\frac{d\sigma}{sT} = (-0.25\pm0.04) \;мН/м\cdot К$$
    $$\epsilon_{\frac{d\sigma}{sT}} = 0.16 \Rightarrow точность $$
    \item Построим требуемые графики: \\
    Пересчет делаем по следующим формулам: 
    $$q = -T \cdot \frac{d\sigma}{d T} - теплота\; образования\; единицы\; поверхности\; жидкости$$ 
    $$\frac{U}{F} = \sigma + q = (\sigma - T \cdot \frac{d\sigma}{d T}) - поверхностная\; энергия\; U\; единицы\; площади$$
    \begin{center}
    \begin{tabular}{|c|c|c|}
    \hline
    $T$, K & $q$, мН & $\frac{U}{F}$, мН\\ \hline
    294    & 74.51 & 138.2          \\ \hline
    298.1  & 75.55 & 138.1          \\ \hline
    303.1  & 76.82 & 138.3          \\ \hline
    308.1  & 78.08 & 138.0          \\ \hline
    314    & 79.58 & 138.0          \\ \hline
    318.1  & 80.62 & 138.1          \\ \hline
    323.1  & 81.88 & 138.0          \\ \hline
    328    & 83.13 & 138.1          \\ \hline
    331.2  & 83.94 & 138.4          \\ \hline
    \end{tabular}
    \end{center}
    \begin{figure}[h]
    \centering
    \includegraphics[width=0.7\linewidth]{2.5.1/два графика.png}
    \label{fig:mpr} 
    \end{figure} 
\end{enumerate}

\section*{Вывод} \\
Первое что стоит заметить, это сходство значений $R$ и $\Delta h$, что бывает не всегда из-за существования "эффективных" радиуса и высоты, в нас же все вычисленные значения попадают в ворота измеренных.\\
В рамках рассматриваемых температур, зависимость $\sigma (T)$ представляется как прямая, причем с хорошей точностью. Сравним полученные значения с табличными:
\begin{center}
\begin{tabular}{|c|c|c|}
\hline
                                    & Табличные значения & Экспериментальные значения \\ \hline
$\sigma_{воды}(20^oC)$, мН/м        & 72.3               & $(63.7\pm2.2)$             \\ \hline
$\frac{d\sigma}{dT}(\approx 30^oC), мН/м\cdot К $ & -0,13              & $(−0.25\pm0.04)$           \\ \hline
$\frac{U}{F}$, мН                   & 115                & 138                        \\ \hline
\end{tabular}
\end{center}\\
В среднем везде наблюдается завышение значений что грустно(




\end{document}